\documentclass[10pt,a4paper]{amsart}
\usepackage[margin=19mm]{geometry}
\usepackage{amsmath,amssymb,mathtools}
\usepackage[scr=rsfs,cal=euler]{mathalfa}
\usepackage{xfrac}
\usepackage[unicode,hidelinks,bookmarks=false]{hyperref}
\newcommand{\ie}{i.e.\ }
\newcommand{\cf}{cf.\ }
\newcommand{\resp}{respectively\ }
\newcommand{\Subsection}{\subsection}
\providecommand{\nobreakdash}{\nobreak\hbox{-}}
\setlength{\emergencystretch}{2em}
\multlinegap0pt
\usepackage{bm}
\usepackage{bbm}
\usepackage{dsfont}
\usepackage{xspace}
\usepackage{cancel}
\usepackage{mathtools}
\usepackage{cleveref}
\usepackage{amsthm}
\makeatletter
\@ifundefined{theorem}{\newtheorem{theorem}{Theorem}}{}
\@ifundefined{lemma}{\newtheorem{lemma}[theorem]{Lemma}}{}
\makeatother
\def\eqref#1{equation~\ref{#1}}
\title[Understanding Generalization in Physics Informed Models thro]{Understanding Generalization in Physics Informed Models through Affine Variety Dimensions}
\author{Takeshi Koshizuka, Issei Sato}
\date{Source: arXiv:2501.18879v2 (CC BY 4.0)}
\begin{document}
\maketitle

\noindent\emph{Source and licence.} Understanding Generalization in Physics Informed Models through Affine Variety Dimensions. Version arXiv:2501.18879v2 (CC BY 4.0), distributed under the Creative Commons Attribution 4.0 International licence, as recorded in the arXiv licence metadata for this exact version and shown on the abstract page https://arxiv.org/abs/2501.18879v2. Full rights notice: licenses/arxiv-2501.18879-CC-BY-4.0.txt.\par\smallskip

\noindent\textit{Editorial scope.} Counted unit: the Lemma whose source label is \texttt{lem:max\_subexp} in the article (source0.tex lines 1723--1732, source label \texttt{lem:max\_subexp}), delivered together with the article's own complete original proof of it (source0.tex lines 1734--1846, one \texttt{proof} environment of 113 lines). No mathematics is written, repaired or re-derived: statement and proof are the article's own text line for line. Required context retained uncounted: none. Every definition, display and label of those lines is retained unchanged. Omitted from this package and not redistributed: 1--1722, 1733--1733, 1847--2358. Nothing inside a retained excerpt is omitted or altered: every retained body is delivered exactly as the article's lines are written, so no omission receipt of any kind accompanies this package. Citations: the retained excerpts carry no citation command at all, so no bibliography entry of the article is transcribed and no reference is redistributed. Reference adaptation: none was needed and none was made; every reference inside the delivered unit resolves inside it, so no unresolved reference or citation is produced. Printed slips kept literal: no printed word, symbol, hypothesis or number of the counted statement or of its proof was altered. Layout and class: the article's own class is \texttt{article} with packages including mystyle, times, amsmath, amsfonts, bm, hyperref, url, authblk, booktabs, nicefrac, microtype, xcolor, amssymb, amsthm; the wrapper is the lane's pinned \texttt{amsart} editorial wrapper at 10pt on A4, which restates the theorem-like environments the retained text needs and adds the article's own macro definitions that the retained text needs (\texttt{\textbackslash eqref}), with extra wrapper packages (bm, bbm, dsfont, xspace, cancel, mathtools, cleveref). The delivered numbering is the unit's own. Assets: no retained span contains a figure environment, a graphics-inclusion command, a PDF-page inclusion, a table or a TikZ picture; no image file is copied into this package, the package declares no asset dependency and no image file is redistributed with it. Licence: version arXiv:2501.18879v2 of this article is distributed under the Creative Commons Attribution 4.0 International licence, as recorded in the arXiv licence metadata for this exact version and shown on the abstract page https://arxiv.org/abs/2501.18879v2. A full rights notice is delivered at licenses/arxiv-2501.18879-CC-BY-4.0.txt. Characters: every retained excerpt is pure ASCII, so the delivered engine, XeLaTeX with its default Latin Modern text font, reports no missing character.
\begin{lemma}[Maximum of sub-exponential random variables]\label{lem:max_subexp}
Let \(X_1,\dots,X_n\) be independent, identically distributed sub-exponential
random variables satisfying \(\|X_i\|_{\psi_1}\le\nu\) for every
\(i \in \{1,\dots,n\}\).  
Then there exists an absolute constant \(C>0\) such that
\begin{align*}
    \left\|\max_{1\le i\le n}\lvert X_i\rvert\right\|_{\psi_1}
    \;\le\; C\,\nu\log n .
\end{align*}
\end{lemma}

\begin{proof}
From the assumption \( \|X_i\|_{\psi_1} \leq \nu \), we have the standard sub-exponential tail bound:
\begin{align*}
    \Pr\left\{ \max_{1 \le i \le n} X_i \ge t \right\}
    &\le
    2\exp\left(
        -\frac{1}{2} \min\left\{
            \frac{t^2}{\nu^2},\;
            \frac{t}{\nu}
        \right\}
    \right).
\end{align*}

Moreover, for all \( t \ge 0 \), we can uniformly bound the maximum via a union bound:
\begin{align}
    \Pr\left\{ \max_{1 \le i \le n} X_i \ge t \right\}
    &= \Pr\left[ \bigcup_{i=1}^{n} \{ X_i \ge t \} \right]
    \le \sum_{i=1}^{n} \Pr\{ X_i \ge t \} \notag \\
    &\le
    2n \exp\left(
        -\min\left\{
            \frac{t^2}{2\nu^{2}},\;
            \frac{t}{2\nu}
        \right\}
    \right).
    \label{eq:union_bound}
\end{align}

Now consider two cases based on the relation between \( n \) and the exponent.

\smallskip
\textbf{Case 1:} Suppose
\begin{align*}
    n \le \exp\left(
        \frac{1}{2} \min\left\{
            \frac{t^2}{2\nu^{2}},\;
            \frac{t}{2\nu}
        \right\}
    \right).
\end{align*}
Then, by inequality \eqref{eq:union_bound}, we have
\begin{align*}
    \Pr\left\{ \max_{i} X_i \ge t \right\}
    \le
    2\exp\left(
        -\frac{1}{2} \min\left\{
            \frac{t^2}{2\nu^2},\;
            \frac{t}{2\nu}
        \right\}
    \right)
    \le
    2\exp\left(
        -\frac{1}{3\log n}
        \min\left\{
            \frac{t^2}{2\nu^2},\;
            \frac{t}{2\nu}
        \right\}
    \right),
\end{align*}
which already provides a stronger tail decay than we ultimately require.

\smallskip
\textbf{Case 2:} Suppose
\begin{align*}
    n > \exp\left(
        \frac{1}{2} \min\left\{
            \frac{t^2}{2\nu^{2}},\;
            \frac{t}{2\nu}
        \right\}
    \right).
\end{align*}
Then, the probability on the right-hand side of \cref{eq:union_bound} satisfies
\begin{align*}
    \Pr\left\{ \max_{1 \leq i \leq n} X_i \ge t \right\}
    \le
    2\exp\left(
        -\frac{1}{3\log n}
        \min\left\{
            \frac{t^2}{2\nu^2},\;
            \frac{t}{2\nu}
        \right\}
    \right)
    > 2e^{-2/3} > 1,
\end{align*}
which is vacuously bounded above by 1. Thus, it does not affect the validity of our bound.

In either case, we obtain the uniform upper bound
\begin{align*}
    \Pr\left\{ \max_{1 \leq i \leq n} X_i \ge t \right\}
    &\le
    2\exp\left(
        -\frac{1}{3\log n}
        \min\left\{
            \frac{t^2}{2\nu^2},\;
            \frac{t}{2\nu}
        \right\}
    \right) \\
    &=
    2\exp\left(
        -\frac{1}{3}
        \min\left\{
            \frac{t^2}{2(\nu\sqrt{\log n})^2},\;
            \frac{t}{2(\nu\log n)}
        \right\}
    \right).
\end{align*}

This tail bound implies that
\begin{align*}
    \left\| \max_{1 \leq i \leq n} X_i \right\|_{\psi_1} \leq C \nu \log n,
\end{align*}
for some universal constant \( C > 0 \).
\end{proof}

\end{document}
